\documentclass[10pt,a4paper]{amsart}
\usepackage[margin=19mm]{geometry}
\usepackage{amsmath,amssymb,mathtools}
\usepackage[scr=rsfs,cal=euler]{mathalfa}
\usepackage{xfrac}
\usepackage[unicode,hidelinks,bookmarks=false]{hyperref}
\newcommand{\ie}{i.e.\ }
\newcommand{\cf}{cf.\ }
\newcommand{\resp}{respectively\ }
\newcommand{\Subsection}{\subsection}
\providecommand{\nobreakdash}{\nobreak\hbox{-}}
\setlength{\emergencystretch}{2em}
\multlinegap0pt
\usepackage{amssymb}
\usepackage{xcolor}
\usepackage{mathtools}
\usepackage{bm}
\usepackage[shortlabels]{enumitem}
\newtheorem{theorem}{Theorem}
\newcommand{\RR}{\mathbb R}
\providecommand{\paragraph}{}
\title{Order-Sensitive Sequential Interventions on Ideal Lattices}
\author{Dmitry Pasechnyuk-Vilensky}
\date{Source: arXiv:2604.26472v1 (CC BY 4.0)}
\begin{document}
\maketitle

\noindent\emph{Source and licence.} Order-Sensitive Sequential Interventions on Ideal Lattices. Version arXiv:2604.26472v1 (CC BY 4.0), distributed by arXiv under the Creative Commons Attribution 4.0 International licence, http://creativecommons.org/licenses/by/4.0/, as recorded in the arXiv licence metadata for this exact version and shown on the abstract page https://arxiv.org/abs/2604.26472v1. Full licence notice: \texttt{licenses/arxiv-2604.26472-CC-BY-4.0.txt}.\par\smallskip

\noindent\textit{Editorial scope.} Counted unit: the article's theorem thm:zero-gauge-app (source0.tex lines 1072--1079). The article's complete original proof of it is delivered with it (source0.tex lines 1081--1189, one \texttt{proof} environment of 109 lines). No mathematics is written, repaired or re-derived: the statement and the proof are the article's own lines, in the article's own notation, and no retained line is substituted or re-flowed.  Retained uncounted as the source-local context the proof needs: lines 243--270.  Nothing was omitted from any retained span, so the package carries no omission record.  No retained line contains a citation, so the wrapper carries no bibliography and no bibliography entry.  No retained line contains a figure environment, an image-inclusion command or drawing code, so no image is delivered and the package declares no asset dependency.  Editorial formatting only, in addition: an AMS \texttt{amsart} preamble that restates the article's own declarations and macros used by the retained lines, namely the environments \texttt{theorem} on separate counters, together with the standard packages those lines need.  The retained environments print in this document as Theorem 1 in order of appearance, whereas the article prints the same items with the numbering of its own full document.  The complete native source of the article as distributed in the arXiv e-print for this exact version is bundled unchanged as \texttt{sources/199396/source0.tex}.
\begin{theorem}[Path-independence criterion]
\label{thm:path-independence-main}
Fix \(x\in\mathcal X\). The following are equivalent:
\begin{enumerate}
    \item for every \(I\subseteq J\) in \(J(P)\) and all \(\gamma,\gamma'\in \Gamma(I,J)\),
    \begin{equation}
    V_x(\gamma)=V_x(\gamma');
    \label{eq:path-independence-condition}
    \end{equation}
    \item for every diamond \(d\),
    \begin{equation}
    \kappa_x(d)=0;
    \label{eq:zero-curvature-condition}
    \end{equation}
    \item there exists a function \(\Phi_x:J(P)\to\RR\) such that for every admissible edge \(I\to I\cup\{a\}\),
    \begin{equation}
    g_x(I,a)=\Phi_x(I\cup\{a\})-\Phi_x(I).
    \label{eq:endpoint-potential-gradient}
    \end{equation}
\end{enumerate}
In this case,
\begin{equation}
V_x(\gamma)=\Phi_x(J)-\Phi_x(I)
\qquad
\forall \gamma\in \Gamma(I,J).
\label{eq:path-independence-telescoping}
\end{equation}
\end{theorem}

\begin{theorem}[Zero-gauge integrability]
\label{thm:zero-gauge-app}
Let \(\kappa\) be a cube-consistent diamond field. Then there exists a unique edge field \(g^{(0)}\) such that
\begin{enumerate}
    \item \(g^{(0)}(K^-,m(K))=0\) for every nonempty ideal \(K\);
    \item the curvature field of \(g^{(0)}\) equals \(\kappa\).
\end{enumerate}
\end{theorem}

\begin{proof}
For an admissible edge \((I,a)\), define
\begin{equation}
\delta(I,a):=\bigl|\{b\in I:\ a<_{\tau} b\}\bigr|.
\label{eq:delta-recursion}
\end{equation}
We construct \(g^{(0)}(I,a)\) recursively on \(\delta(I,a)\).

\paragraph{Base case.}
If \(\delta(I,a)=0\), then \(a\) is the \(<_{\tau}\)-maximum element of \(I\cup\{a\}\), so
\[
(I,a)=((I\cup\{a\})^-,m(I\cup\{a\})).
\]
Define
\begin{equation}
g^{(0)}(I,a):=0.
\label{eq:zero-gauge-base}
\end{equation}

\paragraph{Inductive step.}
Assume \(g^{(0)}\) has been defined on all edges \((K,c)\) with \(\delta(K,c)<r\), and let \((I,a)\) satisfy \(\delta(I,a)=r\ge 1\). Let
\[
b:=m(I),\qquad K:=I\setminus\{b\}.
\]
We first verify that \((K;a,b)\) is a diamond. Since \(a<_{\tau}b\), one has \(a\parallel b\): if \(a\preceq b\), then \(a\in I\) because \(I\) is an ideal, contradiction; if \(b\preceq a\), then \(b<_{\tau}a\) because \(\tau\) extends \(\preceq\), contradiction. Next, \(b\in \mathcal A(K)\) because every predecessor of \(b\) lies in \(I\), and none equals \(b\). Also \(a\in \mathcal A(K)\) because every predecessor of \(a\) lies in \(I\), and \(b\notin \mathrm{Pred}(a)\). Thus \((K;a,b)\) is a diamond.

Now define
\begin{equation}
g^{(0)}(I,a):=g^{(0)}(K,a)+\kappa(K;a,b).
\label{eq:zero-gauge-recursion}
\end{equation}
This is valid because
\[
\delta(K,a)=\delta(I,a)-1.
\]

We now prove by induction on
\[
\Delta(I;u,v):=\bigl|\{b\in I:\ v<_{\tau} b\}\bigr|
\]
that the curvature of \(g^{(0)}\) on every diamond \((I;u,v)\), \(u<_{\tau}v\), equals \(\kappa(I;u,v)\).

\paragraph{Base case.}
Assume \(\Delta(I;u,v)=0\). Then \(v\) is the \(<_{\tau}\)-maximum element of \(I\cup\{u,v\}\), so both edges
\[
I\to I\cup\{v\},
\qquad
I\cup\{u\}\to I\cup\{u,v\}
\]
are reference-tree edges, hence have \(g^{(0)}\)-value \(0\). Therefore
\[
\kappa_{g^{(0)}}(I;u,v)=g^{(0)}(I\cup\{v\},u)-g^{(0)}(I,u).
\]
Applying the recursion \eqref{eq:zero-gauge-recursion} to \((I\cup\{v\},u)\) with \(K=I\) and \(b=v\) gives
\[
g^{(0)}(I\cup\{v\},u)=g^{(0)}(I,u)+\kappa(I;u,v).
\]
Hence \(\kappa_{g^{(0)}}(I;u,v)=\kappa(I;u,v)\).

\paragraph{Inductive step.}
Assume the claim whenever \(\Delta<r\), and let \(\Delta(I;u,v)=r\ge 1\). Set
\[
b:=m(I),\qquad K:=I\setminus\{b\}.
\]
Then \(u,v,b\) are pairwise incomparable. Applying \eqref{eq:zero-gauge-recursion} to the four non-tree edges in the curvature formula gives
\begin{align*}
g^{(0)}(I,v)&=g^{(0)}(K,v)+\kappa(K;v,b),\\
g^{(0)}(I\cup\{v\},u)&=g^{(0)}(K\cup\{v\},u)+\kappa(K\cup\{v\};u,b),\\
g^{(0)}(I,u)&=g^{(0)}(K,u)+\kappa(K;u,b),\\
g^{(0)}(I\cup\{u\},v)&=g^{(0)}(K\cup\{u\},v)+\kappa(K\cup\{u\};v,b).
\end{align*}
Hence
\begin{align*}
\kappa_{g^{(0)}}(I;u,v)
&=
\kappa_{g^{(0)}}(K;u,v)
+\kappa(K;v,b)+\kappa(K\cup\{v\};u,b)-\kappa(K;u,b)-\kappa(K\cup\{u\};v,b).
\end{align*}
By the induction hypothesis,
\[
\kappa_{g^{(0)}}(K;u,v)=\kappa(K;u,v).
\]
Cube consistency on the three-cube \((K;u,v,b)\) gives
\[
\kappa(I;u,v)
=
\kappa(K;u,v)
+\kappa(K;v,b)+\kappa(K\cup\{v\};u,b)-\kappa(K;u,b)-\kappa(K\cup\{u\};v,b).
\]
Comparing the two displays yields
\[
\kappa_{g^{(0)}}(I;u,v)=\kappa(I;u,v).
\]

This proves existence. For uniqueness, let \(\tilde g^{(0)}\) be another such field. Then
\[
h:=g^{(0)}-\tilde g^{(0)}
\]
has zero curvature, so by Theorem~\ref{thm:path-independence-main} there exists \(\varphi:J(P)\to\RR\) such that
\[
h(I,a)=\varphi(I\cup\{a\})-\varphi(I).
\]
Since \(h\) vanishes on every reference-tree edge,
\[
0=h(K^-,m(K))=\varphi(K)-\varphi(K^-)
\qquad \forall K\neq \varnothing.
\]
By induction along the reference tree, \(\varphi\) is constant on all ideals, hence \(h=0\). Therefore \(g^{(0)}=\tilde g^{(0)}\).
\end{proof}

\end{document}
